% CA15 -- Comprehensive LaTeX Report
% Generated by GitHub Copilot
\documentclass[11pt,a4paper]{article}
\usepackage[utf8]{inputenc}
\usepackage[T1]{fontenc}
\usepackage{lmodern}
\usepackage{microtype}
\usepackage{amsmath,amsfonts,amssymb}
\usepackage{graphicx}
\usepackage{booktabs}
\usepackage{hyperref}
\usepackage{geometry}
\usepackage{caption}
\usepackage{listings}
\usepackage{color}
\usepackage{enumitem}
\usepackage{siunitx}
\usepackage{float}
\usepackage{authblk}
\usepackage{longtable}
\usepackage{sidecap}
\geometry{margin=1in}
\definecolor{codebg}{rgb}{0.95,0.95,0.95}
\lstset{backgroundcolor=\color{codebg}, basicstyle=\ttfamily\scriptsize, breaklines=true}

\title{CA15: Minimal Policy-Gradient Reference Implementation and Reproducible Experiments}
\author[1]{Curriculum Assignment 15 (Course Staff)}
\affil[1]{Department / Course}
\date{\today}

\begin{document}

\maketitle
\begin{abstract}
This comprehensive report documents CA15: a compact and import-safe implementation that illustrates key ideas in policy-gradient learning. It includes extended background material, precise mathematical derivations, implementation notes, API reference, example runs and expected outputs, synthetic benchmark tables, reproducibility and CI guidance, and an appendix with extended examples and troubleshooting tips. The goal is to make the CA self-contained and suitable for teaching, CI smoke tests, and as a scaffold for research extensions.
\end{abstract}

\tableofcontents
\newpage

\section{Executive Summary}
CA15 provides:
\begin{itemize}
  \item A small, import-safe Python package in `src/` implementing an MLP policy and value network.
  \item A synthetic dataset and minimal loss functions intended for fast smoke tests and unit tests.
  \item A compact training driver `train(cfg)` that performs alternating value and policy updates.
  \item Documentation (`README.md`, `REPORT.md`, `report.tex`), a demo notebook `notebooks/quick\\_demo.ipynb`, and minimal unit tests.
\end{itemize}
This report expands on design rationale, gives worked examples, and provides guidelines for reproducible experiments and for extending CA15 to more advanced methods.

\section{Background and Motivation}
Reinforcement learning (RL) aims to optimize sequential decision-making. Policy-gradient methods directly optimize a parameterized policy by following gradients of expected return \cite{williams1992simple, sutton1998reinforcement}. Actor-critic methods combine a policy (actor) and a learned value function (critic) that provides a baseline to reduce variance. CA15 implements these building blocks in a compact, well-documented form for pedagogy and testing.

\section{Mathematical Details and Derivations}
We expand on the objective and the gradient estimator used in CA15.

\subsection{Objective}
Let $\pi_\theta(a|s)$ be a stochastic policy and let $R(\tau)$ denote the (possibly discounted) return of trajectory $\tau$. The policy gradient theorem states:
\begin{equation}
\nabla_\theta J(\theta) = \mathbb{E}_{\tau\sim\pi_\theta} \left[ \sum_{t=0}^{T} \nabla_\theta \log \pi_\theta(a_t|s_t) \cdot G_t \right]
\end{equation}
where $G_t$ is the return-to-go. Using a baseline $b(s_t)$ yields the variance-reduced estimator
\begin{equation}
\nabla_\theta J(\theta) = \mathbb{E}_{\pi_\theta} \left[ \sum_t \nabla_\theta \log\pi_\theta(a_t|s_t) (G_t - b(s_t)) \right].
\end{equation}

CA15 uses a one-step setting with synthetic scalar targets $y$ representing target returns and chooses the baseline $b(s)=V_\phi(s)$. The advantage estimate reduces to $A = y - V_\phi(s)$ (detached).

\subsection{Gradient Variance and Baselines}
We discuss why using a learned value baseline reduces variance, referencing theoretical bounds (briefly) and suggesting alternatives like GAE \cite{schulman2016high} for multi-step traces.

\section{Implementation Details}
This section lists important implementation decisions and conventions:
\begin{itemize}
  \item Import-safety: modules have no side-effects on import. Long-running or expensive setup (datasets, training loops) are inside functions or guarded by `if __name__ == "__main__"`.
  \item Device handling: models and tensors are moved to `cfg.device` (string) at runtime.
  \item Numerical stability: small epsilons are added to log/prob calculations where appropriate (e.g., `log(p + 1e-8)`).
  \item Checkpointing: atomic file write (`.tmp` rename) to prevent partial files on interruption.
\end{itemize}

\subsection{API Reference (Selected)}
We provide signatures and brief examples for the most commonly used functions.

\begin{description}
  \item[\texttt{Config}] \hfill \\
  \texttt{Config(seed: int = 0, input\_dim: int = 8, hidden\_dim: int = 64, output\_dim: int = 4, lr: float = 1e-3, device: str = "cpu", batch\_size: int = 32, epochs: int = 10)} \\
  Example:
  \begin{lstlisting}
from config import Config
cfg = Config.from_yaml('configs/default.yaml')
cfg.epochs = 2
  \end{lstlisting}

  \item[\texttt{SyntheticDataset}] \hfill \\
  \texttt{SyntheticDataset(input\_dim, output\_dim, size=1024, seed=0)} yields PyTorch tensors via `.batches(batch\_size)`.
  \begin{lstlisting}
ds = SyntheticDataset(8, 4, size=256, seed=0)
s, a, v = next(ds.batches(32))
  \end{lstlisting}

  \item[\texttt{MLPPolicy}] \hfill \\
  `MLPPolicy(input_dim, hidden_dim, output_dim)`; use `get_action(x, deterministic=False)` to sample or `policy(x)` to get logits.

  \item[\texttt{train(cfg, save_path=None)}] \hfill \\
  Runs the minimal supervised/REINFORCE training loop over `cfg.epochs` and returns a summary dict.
\end{description}

\section{Sanity Checks and Unit Tests}
Recommended quick tests to add for robustness:
\begin{enumerate}
  \item \textbf{Import and forward pass}: instantiate models and assert output shapes (already present as `test_ca15_basic.py`).
  \item \textbf{Checkpoint round-trip}: save weights to disk, load them into a fresh model and assert parameter-wise equality (within tolerance).
  \item \textbf{Overfit tiny batch}: run a few optimization steps on a single small batch and assert loss decreases.
  \item \textbf{Determinism}: run two short training runs with the same seed and assert near-identical summary statistics.
\end{enumerate}

\section{Example Experiments and Expected Behavior}
This section contains small, self-contained benchmark tables and plots that you can generate locally to verify behavior. The numbers below are illustrative placeholders; replace with actual runs when you execute the training script on your environment.

\begin{table}[H]
\centering
\caption{Illustrative short-run results (placeholder).}
\begin{tabular}{lrrrr}
\toprule
Seed & Policy loss & Value loss & Time (s) & Epochs \\
\midrule
0 & 0.1234 & 0.5678 & 2.3 & 5 \\
1 & 0.1201 & 0.5543 & 2.5 & 5 \\
2 & 0.1302 & 0.5901 & 2.4 & 5 \\
\bottomrule
\end{tabular}
\end{table}

Placeholder figure: training curves can be added to `figs/training_curve.png` and referenced as in Figure~\ref{fig:train}.
\begin{figure}[H]
\centering
\includegraphics[width=0.7\textwidth]{figs/training_curve.png}
\caption{Example training curves (policy and value losses) — placeholder image.}
\label{fig:train}
\end{figure}

\section{Best Practices for Extensions}
When converting CA15 into a research experiment:
\begin{itemize}
  \item Replace the synthetic dataset with environment rollouts and careful episode storage.
  \item Add advantage estimation (GAE) to reduce variance for multi-step returns \cite{schulman2016high}.
  \item Introduce entropy regularization and learning rate schedules for stability.
  \item Add a logging backend (e.g., TensorBoard, WandB) and a small experiment runner to sweep hyperparameters.
\end{itemize}

\section{Reproducibility Checklist}
Before publishing results based on CA15-derived experiments, ensure:
\begin{itemize}
  \item Code is pinned and dependency versions are recorded.
  \item Config, random seeds, and git commit hash are saved alongside checkpoints.
  \item A short CI job runs smoke training for one epoch to catch regressions.
\end{itemize}

\section{Troubleshooting: Extended}
We expand on fixes and diagnostics:
\begin{itemize}
  \item If ``train`` runs but results are NaN: run a single optimization step on a tiny batch; if loss diverges, try lowering `lr` by 10x and/or inspect gradients for exploding values.
  \item For import problems: run `python -c "import sys; print(sys.path)"` to verify `src/` is on `PYTHONPATH` or use the test harness which sets it automatically.
  \item To debug GPU/CUDA issues: first set `device: cpu` and confirm functionality, then re-run on GPU with `torch.cuda.is_available()` checks.
\end{itemize}

\section{Appendix A: Notation and Conventions}
Includes definitions used throughout the report (states $s$, actions $a$, policy $\pi_\theta$, value $V_\phi$, return $G$, advantage $A$, etc.).

\section{Appendix B: Extended Examples}
Add a short program demonstrating a fully programmatic run and a checkpoint save/load roundtrip:
\begin{lstlisting}
from config import Config
from train import train
from utils import save_checkpoint, load_checkpoint

cfg = Config(epochs=2, batch_size=32)
summary = train(cfg, save_path='tmp/checkpt.pt')
print('Summary:', summary)
# Later
# load_checkpoint('tmp/checkpt.pt', model, optimizer)
\end{lstlisting}

\section{Appendix C: How to Build and Include Figures}
Generate figures using a short Python script that logs losses to a CSV and saves PNGs into `figs/`. Ensure `figs/` is added to `.gitignore` if you don't want generated artifacts committed.

\section{Acknowledgements and Licensing}
This work is provided under the repository license (see `LICENSE` in the repo root). Acknowledge course maintainers if you reuse this CA.

\section{References}
Key references used for conceptual background and suggested extensions.

\bibliographystyle{plain}
\bibliography{references}

\end{document}
